{\footnotesize
\setlength{\tabcolsep}{3pt}\renewcommand{\arraystretch}{1.22}
\begin{longtable}{@{}p{0.42cm}p{8.3cm}p{4.75cm}@{}}
\toprule
\rowcolor{rpfnavy!12}
\textbf{ที่} & \textbf{ชื่อโครงการ} & \textbf{ตัวชี้วัดที่รับมา}\\
\midrule
\endfirsthead
\toprule
\rowcolor{rpfnavy!12}
\textbf{ที่} & \textbf{ชื่อโครงการ} & \textbf{ตัวชี้วัดที่รับมา}\\
\midrule
\endhead
\multicolumn{3}{@{}l}{\cellcolor{rpfnavy!12}\bfseries ง8-1 ผลักดันเทคโนโลยี นวัตกรรมสู่ชุมชน ตามเป้าหมายการพัฒนาอย่างยั่งยืน}\\
\midrule
1 & โครงการการวิจัย Plant profile ของผักในโรงเรือนอัจฉริยะ & 10:1 $\cdot$ 35:2 $\cdot$ 39:1\\
2 & โครงการการวิจัยพัฒนาวิธีการจัดการกระบวนการผลิตผักเพื่อคุณภาพ & 17:1 $\cdot$ 35:2 $\cdot$ 39:2 $\cdot$ 40:1\\
3 & โครงการพัฒนากระบวนการสกัดน้ำมันอะโวคาโดของศูนย์พัฒนาโครงการหลวงห & 35:7 $\cdot$ 39:2 $\cdot$ 40:1\\
4 & โครงการการพัฒนาตู้อบลมร้อนดอกเก๊กฮวยระบบปั๊มความร้อน เพื่อเพิ่มค & 17:1 $\cdot$ 35:8 $\cdot$ 39:2 $\cdot$ 40:1\\
5 & โครงการพัฒนาเทคโนโลยีเครื่องล้างเมล็ดงาดำแบบกึ่งอัตโนมัติ สำหรับ & 17:1 $\cdot$ 35:12 $\cdot$ 39:2 $\cdot$ 40:1\\
6 & โครงการพัฒนาเทคโนโลยีต้นแบบเครื่องนวดงาขี้ม่อน(งาหอม) เพื่อเพิ่ม & 17:1 $\cdot$ 35:12 $\cdot$ 39:2 $\cdot$ 40:1\\
7 & โครงการประเมินสมบัติเชิงกลของเส้นใยกัญชงจากการอบแห้งด้วยเทคโนโลย & 35:8\\
8 & โครงการศึกษาความเหมาะสมของการนำเศษผักเหลือใช้มาผลิตเป็นถ่านอัดแท & 35:5 $\cdot$ 39:3 $\cdot$ 40:2\\
9 & โครงการ Royal Rose Sensation Giftset ชุดผลิตภัณฑ์ต่อยอดคุณค่ากุห & 17:1 $\cdot$ 35:15 $\cdot$ 38:1 $\cdot$ 39:2 $\cdot$ 40:1\\
10 & โครงการพัฒนากระบวนการผลิตและเครื่องมือการอบแห้งเคพกูสเบอร์รี เพื & 17:1 $\cdot$ 35:3 $\cdot$ 39:2 $\cdot$ 40:1\\
11 & โครงการพัฒนาชุดทดสอบสารพาราควอตเบื้องต้น (Paraquat Test Kit) ที่ & 17:1 $\cdot$ 35:4 $\cdot$ 39:1\\
12 & โครงการนวัตกรรมเชิงพื้นที่ของกระถางเพาะเมล็ดชีวภาพย่อยสลายได้จาก & 35:6 $\cdot$ 39:2 $\cdot$ 40:1\\
13 & โครงการประยุกต์ใช้เทคโนโลยีที่เหมาะสมเพื่อเพิ่มประสิทธิภาพการเลี & 35:6 $\cdot$ 38:1 $\cdot$ 39:3 $\cdot$ 40:2\\
14 & โครงการพัฒนาเครื่องสลัดน้ำผักตามอัตลักษณ์ของศูนย์พัฒนาโครงการหลว & 17:1 $\cdot$ 35:8 $\cdot$ 39:2 $\cdot$ 40:1\\
\midrule
\multicolumn{3}{@{}l}{\cellcolor{rpfnavy!12}\bfseries ง8-2 ขับเคลื่อนกลไกการพัฒนาองค์ความรู้เพื่อยกระดับคุณภาพชีวิต}\\
\midrule
1 & โครงการการถ่ายทอดองค์ความรู้เรื่องการจัดการขยะแบบยั่งยืน ระดับชุ & 35:4\\
2 & โครงการอบรมเชิงปฏิบัติการถ่ายทอดองค์ความรู้เทคนิคการสร้างแม่พิมพ & 35:6 $\cdot$ 39:1 $\cdot$ 40:1\\
3 & โครงการบริหารจัดการโครงการผลิตเมล็ดพันธุ์และปรับปรุงพันธุ์พืชให้ & 35:8 $\cdot$ 36:1\\
4 & โครงการพัฒนาศักยภาพการเรียนรู้ ด้านการจัดการภูมิทัศน์ ผ่านประสบก & 35:7\\
5 & โครงการสนับสนุนการดำเนินงานศูนย์พัฒนาพันธุ์พืชจักรพันธ์เพ็ญศิริ & 35:5 $\cdot$ 36:1 $\cdot$ 39:3 $\cdot$ 40:3\\
6 & โครงการพัฒนาแหล่งเรียนรู้เกษตรอันเนื่องมาจากพระราชดำริ วิถีธรรมช & 35:7 $\cdot$ 35:5 $\cdot$ 36:1 $\cdot$ 39:3 $\cdot$ 39:1 $\cdot$ 40:1 $\cdot$ 40:1\\
7 & โครงการจัดทำคำขอสิ่งบ่งชี้ทางภูมิศาสตร์ไทย (GI) สินค้ากาแฟเลอตอ & 38:1 $\cdot$ 39:1 $\cdot$ 40:1\\
8 & โครงการพัฒนาพื้นที่ต้นแบบ โคก หนอง นา เพื่อการเรียนรู้ด้านเกษตรอ & 10:1 $\cdot$ 36:1 $\cdot$ 39:1\\
9 & โครงการประยุกต์ใช้การคำนวนสูตรอาหารไก่ไข่ต้นทุนต่ำด้วยโปรแกรมสำเ & 35:4 $\cdot$ 39:1 $\cdot$ 40:1\\
10 & โครงการถ่ายทอดองค์ความรู้และนวัตกรรมการยกระดับวัสดุเพาะเห็ดเศรษฐ & 35:8 $\cdot$ 39:1 $\cdot$ 40:1\\
11 & โครงการเพิ่มมูลค่าสินค้าเกษตรและการพัฒนาสื่อด้วย AI พื้นที่เขตปฏ & 35:4 $\cdot$ 39:1 $\cdot$ 40:1\\
\midrule
\multicolumn{3}{@{}l}{\cellcolor{rpfnavy!12}\bfseries ง8-3 พัฒนากำลังคน สร้างอาชีพ ลดความเหลื่อมล้ำ บนพื้นที่สูง}\\
\midrule
1 & โครงการพัฒนาฐานข้อมูลด้านวิศวกรรมและยกระดับทักษะอาชีพในพื้นที่โค & 35:1 $\cdot$ 39:1 $\cdot$ 40:1\\
2 & โครงการยกระดับการบริหารจัดการ ระบบติดตามและเครือข่าย ความร่วมมือ & 10:5\\
3 & โครงการการจัดกิจกรรมการเรียนรู้ส่งเสริมงานวิจัยและบริการวิชาการ & 35:4\\
4 & โครงการการฝึกอบรมสร้างความสามารถการออกแบบระบบจ่ายน้ำในแปลงเกษตร & 35:1 $\cdot$ 39:1 $\cdot$ 40:2\\
5 & โครงการพัฒนาทักษะอาชีพด้านพลังงานสะอาด & 35:3 $\cdot$ 39:1 $\cdot$ 40:1\\
6 & โครงการนำร่อง SkillChain มทร.ล้านนา :ระบบรับรองทักษะช่างและจัดหา & 10:2\\
7 & โครงการการสำรวจผลผลิตทางการเกษตรโครงการหลวง เพื่อการประกอบอาหารแ & 35:8\\
8 & โครงการยกระดับมาตรฐานผลิตภัณฑ์ผ้าใยกัญชงวิสาหกิจชุมชนดาวม่างสู่ & 35:5 $\cdot$ 38:1 $\cdot$ 39:2 $\cdot$ 40:2\\
9 & โครงการยกระดับห่วงโซ่คุณค่าห้อมพื้นถิ่น สู่ผลิตภัณฑ์ย้อมสี ธรรมช & 35:9 $\cdot$ 38:1 $\cdot$ 39:2 $\cdot$ 40:1\\
10 & โครงการการพัฒนานวัตกรรมด้านการศึกษาเพื่อการพัฒนาศักยภาพเยาวชนในโ & 35:6 $\cdot$ 39:1 $\cdot$ 40:1\\
11 & โครงการประเมินผลกระทบของโครงการบริการวิชาการมหาวิทยาลัยเทคโนโลยี & 35:7\\
12 & โครงการพัฒนาแผนการตลาดของกลุ่มวิสาหกิจชุมชนภายใต้แนวคิด “กินได้ & 35:7 $\cdot$ 38:2 $\cdot$ 39:1\\
13 & โครงการต้นแบบอาคารอนุรักษ์พลังงานในรูปแบบ (Green Building \& Gree & 35:3\\
14 & โครงการพัฒนาทักษะวิชาชีพให้กับเจ้าหน้าที่ ชาวบ้าน เยาวชน พื้นที่ & 35:7 $\cdot$ 35:7 $\cdot$ 39:1 $\cdot$ 39:1 $\cdot$ 40:1 $\cdot$ 40:1\\
15 & โครงการพัฒนาต้นแบบชุดน้ำชาเซรามิกสโตนแวร์สำหรับยกระดับผลิตภัณฑ์ช & 35:8\\
16 & โครงการการถ่ายทอดองค์ความรู้กระบวนการผลิตน้ำดื่มสะอาดด้วยระบบรถน & 35:5 $\cdot$ 39:1 $\cdot$ 40:1\\
17 & โครงการพัฒนาทักษะวิชาชีพให้กับเจ้าหน้าที่ ชาวบ้าน เยาวชน พื้นที่ & 35:7 $\cdot$ 35:7 $\cdot$ 39:1 $\cdot$ 39:1 $\cdot$ 40:1 $\cdot$ 40:1\\
18 & การประยุกต์ใช้ lora mesh ชนิด multi-hops ในการส่งสัญญาณจากเซนเซอ & 35:3 $\cdot$ 39:1 $\cdot$ 40:1\\
19 & การประยุกต์ใช้เครือข่าย lora-mesh ชนิด multi-hops เพื่อการส่งสัญ & 35:3 $\cdot$ 39:1 $\cdot$ 40:1\\
20 & โครงการวิจัยความเข้มแข็งทางสังคม ชุมชนในการวางแผนการใช้น้ำศูนย์พ & 35:1\\
21 & โครงการพัฒนาทักษะอาชีพและองค์ความรู้ด้านเกษตรและเทคโนโลยีอย่างยั & 35:5 $\cdot$ 39:3 $\cdot$ 40:3\\
22 & โครงการถ่ายทอดเทคโนโลยีการเพาะขยายพันธุ์ปลาเลียหินเพื่อการอนุรัก & 35:3 $\cdot$ 39:1 $\cdot$ 40:1\\
23 & โครงส่งเสริมพัฒนาโจทย์วิจัยจากชุมชน & 35:7\\
24 & โครงการถ่ายทอดองค์ความรู้และนวัตกรรมเพื่อยกระดับมูลค่าไก่ดำ (ไก่ & 35:12 $\cdot$ 39:1 $\cdot$ 40:1\\
25 & โครงการพัฒนาระบบติดตามข้อมูลภายใต้สถานีตรวจวัดสภาพอากาศเพื่อป้อง & 35:3\\
26 & โครงการจัดทำคำขอสิ่งบ่งชี้ทางภูมิศาสตร์ไทย (GI) สินค้ากาแฟเลอตอ & 38:1 $\cdot$ 39:1 $\cdot$ 39:1 $\cdot$ 40:1 $\cdot$ 40:1\\
\midrule
\bottomrule
\end{longtable}
}
